\section{Introduction}
\label{sec:intro}

Recent advances in video diffusion models demonstrate significant improvements in visual and motion quality ~\citep{VideoDiffusionModels,blattmann2023stable,AlignYourLatents,LVDM}. 
These models typically take a text prompt~\citep{VideoDiffusionModels,blattmann2023stable,ImagenVideo} or image~\citep{VideoCrafter1,VideoCrafter2,AnimateDiff,DynamiCrafter} as input and generate video frames of a photorealistic, animated scene.
Current methods can generate videos that are largely consistent with an input text description or image; however, fine-grained adjustment of specific video elements (e.g., object motion or camera movement) is conventionally a tedious process that requires re-running the model with different text prompts or random seeds~\citep{wu2023freeinit,qiu2024freetraj}.

Approaches for controllable video generation aim to eliminate this process of trial-and-error through direct manipulation of generated video elements, such as object motion~\citep{wu2024draganything,yin2023dragnuwa,wang2024boximator}, pose~\citep{AnimateAnyone,MagicAnimate}, and camera movement~\citep{MotionCtrl,li2024image,CameraCtrl,MotionMaster}.
One line of work fine-tunes pre-trained video generators to incorporate control signals such as bounding boxes or point trajectories~\citep{wu2024draganything,MotionCtrl}. 
One of the primary challenges with these supervised methods is the expensive training cost, and thus, previous methods usually incorporate trajectory control by fine-tuning at a lower resolution than the original model~\cite{wu2024draganything, yin2023dragnuwa}.
More recently, several methods for zero-shot, controllable text-to-video generation have been developed~\citep{ma2023trailblazer,qiu2024freetraj,jain2024peekaboo}.
They control object trajectories by modulating the cross-attention maps between features within a bounding box and an object-related text token.
Still, it is not always possible to associate a desired edit with the input text prompt (consider, e.g., motion of object parts). 
Moreover, these methods cannot be directly applied to animate existing images, as they are only conditioned on text.

In this work, we propose SG-I2V, a new method for controllable image-to-video generation.
Our approach is \textit{self-guided}, in that it offers zero-shot control by relying solely on knowledge present in a pre-trained video diffusion model. 
Concretely, given an input image, a user specifies a set of bounding boxes and associated trajectories.
Then, our framework alters the generation process to control the motion of target scene elements.
It is essential to manipulate the structure of the
generated video to achieve precise control over element positions, which is mainly decided by early denoising steps~\citep{balaji2022ediff,DMProcessAnalyze}.
In image diffusion models, it is known that feature maps extracted from the output of upsampling blocks are \textit{semantically aligned}, i.e., pixels belonging to the same object share similar feature vectors on the feature map and thus can be used to control the spatial layout of generated images~\citep{tang2023emergent, shi2024dragdiffusion, namekata2024emerdiff, tumanyan2023plug}.
However, our analysis reveals that feature maps extracted from the upsampling blocks of video diffusion models are only weakly aligned across frames (see~\cref{fig:correspondence}).
This misalignment poses challenges, as directly manipulating these feature maps fails to give useful guidance signals for layout control. %
Instead, we find that feature maps extracted from the self-attention layers can be semantically aligned by replacing the key and value tokens for each frame with those of the first frame (see bottom row of \cref{fig:correspondence}).
After that, we can control the motion of generated videos by optimizing the latent (the input to the denoising network) with a loss that encourages similarity between the aligned features within each bounding box along the input trajectory. 
Finally, we apply a post-processing step to enhance output quality by ensuring that our optimization does not disrupt the distribution of high-frequency noise expected by the diffusion model.




In summary, our work makes the following contributions:
\begin{itemize}
    \item We conduct a first-of-its-kind analysis of semantic feature alignment in a pre-trained image-to-video diffusion model and identify important differences from image diffusion models.
    \item Building on this analysis, we propose SG-I2V, a zero-shot, self-guided approach for controllable image-to-video generation.
    Our method can control object motion and camera dynamics for arbitrary input images and any number of objects or regions of a scene.
    \item We conduct extensive experiments to show superior performance over zero-shot baselines while significantly narrowing down the performance gap with supervised baselines in visual and motion quality.
\end{itemize}





\begin{figure}[!t]
\centering
\includegraphics[width=\textwidth]{imgs/semantic_correspondence_labeled.jpg}
\caption{\textbf{Semantic correspondences in video diffusion models.}
We analyze feature maps in the image-to-video diffusion model SVD~\citep{blattmann2023stable} for three generated video sequences (row 1). 
We use PCA to visualize the features at diffusion timestep 30 (out of 50) at the output of an upsampling block (row 2), a self-attention layer (row 3), and the same self-attention layer after our alignment procedure (row 4). %
Although output feature maps of upsampling blocks in image diffusion models are known to encode semantic information ~\citep{tang2023emergent}, we only observe weak semantic correspondences across frames in SVD. 
Thus, we focus on the self-attention layer and modify it to produce feature maps that are semantically aligned across frames.
}
\label{fig:correspondence}
\vspace{-2.0mm}
\end{figure}



